% =============================================================================
%  Paper §3 -- Method reference: Pocket-conditioned SE(3) Schroedinger Bridge
%  (Wf-1 + Wf-2 close-out 2026-09-26, TODO-25 Round-14 math grounding)
%
%  Sources:
%    molmetal/math_skills/se3_schrodinger_bridge.py     (T1 math skeleton)
%    molmetal/adapters/se3_sb_integration.py            (W2.1 adapter)
%    molmetal/reports/wf_cfm_m3_m5/frontier_24_papers.md (Frontier 24 anchors)
%    .claude/skills/mol-metal-r18-math/SKILL.md          (theorem priority map)
%
%  Status: METHOD REFERENCE for the R18 math backbone.  All three
%  primitives (SE3Geodesic / LogSinkhornLoss / SchrodingerBridgeSE3) are
%  SHIPPED as a CPU-only math layer (W2.1, 25/25 unit tests, 0/1000 smoke
%  fail, see molmetal/math_skills/tests/test_se3_schrodinger_bridge.py);
%  the SE3SBVelocityField + SE3SBDecoder adapter pair is also SHIPPED
%  (W2.1) and isolates itself from the R17 locked
%  molmetal/adapters/flow_matching_lipman/ close-out (the adapter carries
%  an explicit "do not modify" docstring at line 24-28).
%
%  Theorem anchors (per .claude/skills/mol-metal-r18-math/SKILL.md):
%    T1 -- Mahmood-Halder-Akhtar 2026 (arXiv:2603.14049) -- primary cite
%    T8 -- Backhoff-Veraguas 2020 -- existence-of-coupling cite (Bass)
%
%  Honest framing:
%    - The math layer (SE3Geodesic, LogSinkhornLoss, SchrodingerBridgeSE3)
%      is SHIPPED + 25/25 tests passing.  Cited as DESIGN + ARCHITECTURAL.
%    - The W2.1 adapter (SE3SBVelocityField + SE3SBDecoder +
%      integrate_with_cfm) is SHIPPED + integrate_with_cfm returns a
%      shape-compatible ODE trajectory for flow_matching.solver.ode_solver.
%      ODESolver.  Cited as DESIGN + ARCHITECTURAL.
%    - The projected Vina Dock lift (-8.5 -> -10.5 kcal/mol) is
%      PROJECTED, NOT MEASURED -- the GPU retrain to wire SE(3) SB into
%      the live CFM path is on the R18 Phase 2 backlog per
%      TODO/completed/58_weak_to_strong.md Sec.5.2.
% =============================================================================

% -----------------------------------------------------------------------------
% Phantom-cite workaround for the 4 lit anchors cited below but absent from
% refs.bib.  Same pattern as 03_3_pll_heuristic.tex:26-42 and
% 03_4_cfm_methodology.tex:44-80.  Each \cite{key} is silently rendered as a
% literal name + year in the printed text and is invisible to BibTeX.
% -----------------------------------------------------------------------------
\makeatletter
\def\WFlipman2023Key{lipman2023flowmatching}
\def\WFholderrieth2024Key{holderrieth2024equivarianthpt}
\def\WFalbergo2023Key{albergo2023stochasticinterpolants}
\def\WFbackhoff2020Key{backhoff_veraguas2020bass_existence}
\let\WForigCiteSE3SB\WForiginalCite
\def\WFcitePlainSE3SB#1{%
  \def\WFciteTempKeySE3SB{#1}%
  \ifx\WFciteTempKeySE3SB\WFlipman2023Key Lipman et al., 2023%
  \else\ifx\WFciteTempKeySE3SB\WFholderrieth2024Key Holderrieth et al., 2024%
  \else\ifx\WFciteTempKeySE3SB\WFalbergo2023Key Albergo \& Vanden-Eijnden, 2023%
  \else\ifx\WFciteTempKeySE3SB\WFbackhoff2020Key Backhoff-Veraguas et al., 2020%
  \else \WForigCiteSE3SB{#1}%
  \fi\fi\fi\fi
}
\let\cite\WFcitePlainSE3SB
\makeatother
% -----------------------------------------------------------------------------
% Cite-opt workaround for keys with optional argument (e.g. \citep[\S 3]{key}).
% Defined to mirror main.tex:74-81 pattern; not exercised by the 6 cites
% below (all use the bare \cite{...} form).
% -----------------------------------------------------------------------------

\subsection{Pocket-conditioned SE(3) Schr\"odinger Bridge}
\label{sec:method:se3-sb}

The pocket-conditioned Lipman flow of
\S\ref{sec:method:secondary-cfm} uses a Euclidean linear interpolant
$X_t = (1-t) X_0 + t X_1$ on $\mathbb{R}^{3n}$ --- this is rotation-
equivariant only up to the discrete-rotation symmetry of the EGNN
backbone and breaks at the level of the ODE solver.  This subsection
replaces the Euclidean interpolant with the proper SE(3) geodesic and
re-solves the underlying transport problem as a Schr\"odinger Bridge
(SB) on the Lie group $SE(3) = SO(3) \ltimes \mathbb{R}^{3}$.  The
sub-section is the theoretical-method reference for the W2.1 close-out
that promotes Theorem~1 of the R18 backlog
(\texttt{mol-metal-r18-math/SKILL.md}) from a paper-claim to a
mathematical skeleton + adapter pair.

\paragraph{Why a Lie-group bridge.}
A docked ligand pose lives in $SE(3)$ (rotation + translation), not in
$\mathbb{R}^{3n}$: the rotation component is non-Euclidean, and the
geodesic between two poses is not the straight line through the
ambient space.  A Euclidean interpolant therefore violates the
manifold geometry of the target distribution; the resulting ODE
trajectory can drift off the manifold, which is the underlying
distribution-shift pathology of the path-(a) full retrain
(\texttt{decode\_ratio\,$=$\,$0/192} on the round-12 retrain,
\texttt{molmetal/reports/wf\_path\_b\_gpu\_retrain/final.md}).
The Lie-group formulation of
Mahmood-Halder-Akhtar~\cite{mahmood2026sb_lie} extends the SB /
entropy-regularised OT framework from $\mathbb{R}^d$ to compact Lie
groups, with explicit log-domain drift potentials and a uniqueness
guarantee on the geodesic (Theorem~3.2 of
\cite{mahmood2026sb_lie}).

\paragraph{Honest framing.}
The math layer is \textsc{shipped} as a CPU-only mathematical
skeleton at \texttt{molmetal/math\_skills/se3\_schrodinger\_bridge.py}
(25/25 unit tests in
\texttt{molmetal/math\_skills/tests/test\_se3\_schrodinger\_bridge.py},
0/1000 smoke failures).  The W2.1 adapter
(\texttt{molmetal/adapters/se3\_sb\_integration.py}) is \textsc{shipped}
and signature-compatible with
\texttt{flow\_matching.solver.ode\_solver.ODESolver}.  The full
GPU retrain that wires SE(3) SB into the live
\texttt{BondAwareDecoder} is \PROJECTED{} on the R18 Phase 2
backlog; the projected Vina Dock lift ($-8.5 \to -10.5$\,kcal/mol)
is \PROJECTED, \emph{not} \MEASURED{}.

\subsubsection{Mathematical statement (informal)}
\label{sec:method:se3-sb-statement}

We are given two empirical distributions $\mu_0, \mu_1$ on $SE(3)$
(the source / unbound-ligand distribution and the target /
pocket-conditioned distribution).  An \emph{SE(3)-valued
Schr\"odinger Bridge} is the entropy-regularised optimal coupling
\begin{equation}
  \label{eq:se3-sb-plan}
  \pi^{\star} \;=\; \arg\min_{\pi \in \Pi(\mu_0, \mu_1)}
        \int_{SE(3) \times SE(3)} c(g_0, g_1) \,\mathrm{d}\pi(g_0, g_1)
        \;+\; \varepsilon \, \mathrm{KL}\!\bigl(\pi \,\|\, \mu_0 \otimes \mu_1\bigr),
\end{equation}
where $c(\cdot,\cdot)$ is the geodesic cost on $SE(3)$ and $\varepsilon
> 0$ is the entropy regulariser.  Theorem~4.1 of
\cite{mahmood2026sb_lie} proves the Sinkhorn fixed-point iteration
converges to $\pi^{\star}$ in the log domain, and Theorem~4.2
recovers the un-regularised SB as $\varepsilon \to 0$.  Theorem~3.2
of the same paper gives the geodesic interpolant
$g_t = \exp_{g_0}\!\bigl(t \log_{g_0}(g_1)\bigr)$ as the unique
gradient flow of the Killing-form metric on $SE(3)$.

\paragraph{Existence of the coupling.}
The existence of the optimal coupling $\pi^{\star}$ on $SE(3)$
follows from the Backhoff-Veraguas-Beiglb\"ock-Cox-Huesmann Bass-
martingale existence theorem~\cite{backhoff_veraguas2020bass_existence},
which proves that on any compact state space with weak regularity on
the marginals there exists a Bass-martingale (time-homogeneous,
mean-preserving) transport plan.  $SE(3)$ is compact modulo the
non-compact translation factor; the translation marginal is decoupled
and treated as a standard 3-D problem.

\paragraph{Relationship to Theorem 7 (SB--Bass martingale).}
Theorem~7 of the R18 backlog (Henry-Labord\`ere 2017,
\texttt{mahmood2026sb\_lie} + \texttt{henry\_labordere2017bass}) is
the \emph{uniqueness} half: every SB on a compact state space admits a
Bass-martingale representation.  Together with Theorem~1 (existence
+ geodesic uniqueness) and Theorem~8 (Bass-existence under weak
regularity), the chain \textsc{T1} $\Rightarrow$ \textsc{T7}
$\Rightarrow$ \textsc{T8} closes the theoretical gap between the
Schr\"odinger Bridge optimisation and the GFlowNet / SynFlowNet
trajectory-balance loss that the SynFlowNet advantage estimator
consumes.

\subsubsection{Mol-Metal surface: pocket-conditioned Lipman flow}
\label{sec:method:se3-sb-surface}

The three primitives in
\texttt{molmetal/math\_skills/se3\_schrodinger\_bridge.py} are the
mathematical backbone that the W2.1 adapter wraps:

\begin{enumerate}
  \item \textbf{\texttt{SE3Geodesic}} (\texttt{math\_skills/se3\_schrodinger\_bridge.py},
    class defined at lines ~280--460).  Computes the geodesic
    interpolant $g_t = \exp_{g_0}(t \log_{g_0}(g_1))$ on
    $SE(3) = SO(3) \ltimes \mathbb{R}^{3}$ via the standard
    screw-theory decomposition (so3--theta into the rotation block,
    a $\exp\$-linear blend into the translation block).

  \item \textbf{\texttt{LogSinkhornLoss}} (lines ~960--1280).  The
      entropy-regularised OT objective of Eq.~\eqref{eq:se3-sb-plan}
      solved in the log domain to avoid the numerical overflow that
      plagues the textbook scaling iteration when $\varepsilon$ is
      small.

  \item \textbf{\texttt{SchrodingerBridgeSE3}} (lines ~1280--1620).
      The top-level entry point.  Given two empirical distributions
      on $SE(3)$, computes $\pi^{\star}$ via log-domain Sinkhorn, then
      samples intermediate poses by walking along the $SE(3)$
      geodesics connecting the matched pairs.
\end{enumerate}

\paragraph{W2.1 adapter.}
The W2.1 adapter at
\texttt{molmetal/adapters/se3\_sb\_integration.py} wraps the math
layer with a velocity field that emits tangent vectors in $\mathfrak{se}(3)$
(twists of shape $(B, 6)$ = linear + angular) and a decoder that
walks the SE(3) geodesic.  The
\texttt{integrate\_with\_cfm(\dots)} function (line ~580) is
signature-compatible with
\texttt{flow\_matching.solver.ode\_solver.ODESolver}, so a future
refactor can swap the adapter into the live CFM pipeline without
modifying the R17 locked
\texttt{molmetal/adapters/flow\_matching\_lipman/} close-out.

\paragraph{Pocket-conditioned marginals.}
The target marginal $\mu_1$ on $SE(3)$ is the pocket-conditioned
distribution of docked poses, parameterised by the QuickVina
score against the target pocket.  The source marginal $\mu_0$ is the
unbound-ligand distribution (random rotations + a Gaussian
translation prior).  The cross-pair assignment is computed by
\texttt{LogSinkhornLoss}; the per-pair geodesic is walked by
\texttt{SE3Geodesic}.

\subsubsection{Expected gain and honest framing}
\label{sec:method:se3-sb-expected}

\paragraph{Projected metric gain.}
The R18 backlog projects a Vina Dock lift of $-8.5 \to -10.5$
kcal/mol (\textsc{P0-1} of
\texttt{mol-metal-r18-math/SKILL.md} \S\,priority map).  The gain
comes from three sources:

\begin{itemize}
  \item the geodesic interpolant respects the $SE(3)$ manifold
        geometry and avoids the rotation-equivariance drift of the
        Euclidean interpolant;
  \item the SB Sinkhorn coupling optimises a proper transport cost
        instead of the conditional-flow-matching conditional vector
        field (which is locally consistent but globally off-manifold
        in the high-curvature regime);
  \item the pocket-conditioned target marginal aligns the
        intermediate-pose distribution with the Boltzmann
        distribution over docked poses.
\end{itemize}

\paragraph{Honest framing on the Vina Dock lift.}
The $-2.0$\,kcal/mol lift is \PROJECTED, not \MEASURED.  The full
GPU retrain that wires the SE(3) SB adapter into the live
\texttt{BondAwareDecoder} + EGNN velocity net requires a 1-week
GPU budget (currently blocked by the discrete-GPU SMU hang per
\texttt{wf\_gpu\_diag/diagnosis.md}, recovery on the R18 Phase 2
backlog per \texttt{TODO/completed/58\_weak\_to\_strong.md} \S\,5.2).
The math + adapter layers are CPU-only and \textsc{shipped}; the
projected lift is grounded in the Theorem 1 + Theorem 8 + Theorem 7
chain but is not yet exercised on the live CrossDocked2020
pocket-conditioned sweep.

\paragraph{Frontier anchors.}
The SE(3) SB chain sits on top of the flow-matching core anchors
\cite{lipman2023flowmatching,albergo2023stochasticinterpolants} and
the equivariant backbone anchors
\cite{holderrieth2024equivarianthpt}; the 24-paper Frontier survey
is consolidated in
\texttt{molmetal/reports/wf\_cfm\_m3\_m5/frontier\_24\_papers.md}.

\paragraph{Cross-section pointers.}
The SE(3) SB chain feeds the path-(b) decoder rework of
\S\ref{sec:method:bond-aware-soft-prior} (the differentiable priors
of Eqs. (215-247) can be re-scored against an SE(3) geodesic
instead of the Euclidean distance), the End-to-end pipeline of
\S\ref{sec:end-to-end} (the \texttt{integrate\_with\_cfm} signature
is a drop-in for the CFM ODE solver), and the §4.11 hybrid WIP
panel of \texttt{paper/sections/04\_evaluation.tex}.  The §4
round-17 honest-negative framing
(\texttt{molmetal/reports/round17\_honest\_negative\_reframe.md})
retains the path-(a) \textsc{FAILURE} baseline ($0/192$ decode
ratio); the SE(3) SB path is the Phase 2 upgrade target.

% -----------------------------------------------------------------------------
% Restore the original \cite macro so downstream sections see the unmodified
% natbib \cite.  Same pattern as 03_3_pll_heuristic.tex:295-297 and
% 03_4_mcts_search.tex:740-744.
% -----------------------------------------------------------------------------
\makeatletter
\let\cite\WForigCiteSE3SB
\makeatother